%!TEX root = ../thesis.tex

\abstractUkr

    Кваліфікаційна робота містить: \pageref{LastPage}~стор., \totfig~рисунків, \tottab~таблиць, \totbib~джерел.

    Мета роботи -- розробка обчислювально ефективної системи інтерактивної багатоваріантної колоризації.
    Об'єкт дослідження -- процес керованого відновлення кольору на монохромних зображеннях.
    Предмет дослідження -- гібридні генеративні нейромережеві моделі на базі згорток та селективних моделей простору станів.

    Спроєктовано та реалізовано гібридну архітектуру B-Mamba, що поєднує точність локального вилучення ознак U-Net із глобальним рецептивним полем двонаправленої Mamba.
    Для багатоваріантності розроблено алгоритм стохастичної генерації підказок на основі градієнтної значущості.
    Впровадження спільних ваг дозволило мінімізувати апаратні вимоги.
    Створено вебзастосунок для керування колоризацією у реальному часі.

    Експериментальний аналіз показав, що модель перевершує базові моделі, досягнувши FID 3.078 та LPIPS 0.143.
    Забезпечено низьке значення похибки дотримання підказок ($\mathrm{MSE}_{\mathrm{HINT}} = 6.1 \times 10^{-5}$) при інференсі 15.08 мс та споживанні 187 МБ відеопам'яті.
    Користувацьке дослідження підтвердило високу реалістичність генерації та точність просторового керування кольором.

    \MakeUppercase{КОЛОРИЗАЦІЯ ЗОБРАЖЕНЬ, ГЛИБОКЕ НАВЧАННЯ, U-NET, MAMBA, СЕЛЕКТИВНІ МОДЕЛІ ПРОСТОРУ СТАНІВ, ІНТЕРАКТИВНА КОЛОРИЗАЦІЯ, БАГАТОВАРІАНТНА КОЛОРИЗАЦІЯ}


\abstractEng

    The qualification work contains: \pageref{LastPage}~pages, \totfig~figures, \tottab~tables, \totbib~references.

    The purpose of the work is to develop a computationally efficient system for interactive multi-variant image colorization.
    The object of research is the process of controlled color restoration in monochrome images.
    The subject of research is hybrid generative neural network models based on convolutions and selective state space models.

    A hybrid architecture, B-Mamba, was designed and implemented, combining the precision of U-Net local feature extraction with the global receptive field of bidirectional Mamba.
    For multi-variant colorization, an algorithm for stochastic hint generation based on gradient saliency was developed.
    The introduction of shared weights minimized hardware requirements.
    A web application was created for real-time colorization control.

    Experimental analysis showed that the model outperforms baseline models, achieving an FID of 3.078 and LPIPS of 0.143.
    A low value of hint-following error was achieved ($\mathrm{MSE}_{\mathrm{HINT}} = 6.1 \times 10^{-5}$) with a 15.08 ms inference time and 187 MB of video memory consumption.
    A user study confirmed the high visual realism of the generation and the accuracy of spatial color control.

    \MakeUppercase{IMAGE COLORIZATION, DEEP LEARNING, U-NET, MAMBA, SELECTIVE STATE SPACE MODELS, INTERACTIVE COLORIZATION, MULTI-VARIANT COLORIZATION}

\clearpage
% %!TEX root = ../thesis.tex

% \abstractUkr

% Кваліфікаційна робота містить: ??? стор., ??? рисунки, ??? таблиць, ??? джерел.

% У рефераті роботи ви повинні коротко (два-три абзаци) викласти, що саме 
% було зроблено у цій роботі. Перші три речення реферату (після статистичних 
% даних) повинні окреслити мету роботи, об'єкт та предмет дослідження. Після 
% цього викладаються основні результати, одержані в ході дослідження.

% Наприкінці анотації великими літерами зазначаються ключові слова. Ось так:

% \MakeUppercase{КЛЮЧОВІ СЛОВА, СИМЕТРИЧНА КРИПТОГРАФІЯ, ФІЗТЕХ НАЙКРАЩІЙ}


% \abstractEng

% The English abstract must be the exact translation of the Ukrainian 
% ``annotation'' (including statistical data and keywords).

% \clearpage